\section{Processes}
A \textcolor{Blue}{\textbf{process}} is a running instance of a \textcolor{Bittersweet}{program} (static, on disk). It bundles three contexts:
\begin{itemize}
  \item \textcolor{Bittersweet}{Memory}: \textcolor{Bittersweet}{text} (code, RO), \textcolor{Bittersweet}{data} (globals/statics), \textcolor{Bittersweet}{heap} (grows up), \textcolor{Bittersweet}{stack} (grows down).
  \item \textcolor{Bittersweet}{Hardware}: GPRs, \textcolor{Bittersweet}{PC}, \textcolor{Bittersweet}{SP}, status flags.
  \item \textcolor{Bittersweet}{OS} (\textcolor{Blue}{PCB}): PID, PPID, state, saved regs+PC, MMU info, FD table, scheduling info, signal handlers, accounting.
\end{itemize}

\subsection{Process states}
\begin{center}
\begin{tikzpicture}[every node/.style={font=\scriptsize}, scale=0.9, ->,>=Stealth,
  state/.style={draw, rounded corners, inner sep=2pt, font=\scriptsize}]
  \node[state] (N) at (0,1.2) {New};
  \node[state] (R) at (2.0,1.2) {Ready};
  \node[state] (Run) at (4.4,1.2) {Running};
  \node[state] (B) at (4.4,-0.2) {Blocked};
  \node[state] (T) at (6.6,1.2) {Terminated};
  \draw (N) -- node[above,font=\tiny] {admit} (R);
  \draw ([yshift=2pt]R.east) -- node[above,font=\tiny] {dispatch} ([yshift=2pt]Run.west);
  \draw ([yshift=-2pt]Run.west) -- node[below,font=\tiny] {preempt} ([yshift=-2pt]R.east);
  \draw (Run) -- node[right,font=\tiny] {wait} (B);
  \draw (B) -- node[below,font=\tiny] {event} (R);
  \draw (Run) -- node[above,font=\tiny] {exit} (T);
\end{tikzpicture}
\end{center}
\textcolor{ForestGreen}{Unix variant} adds \textcolor{ForestGreen}{Stopped} (suspended via \texttt{SIGSTOP}; resumed via \texttt{SIGCONT}).
\textcolor{Bittersweet}{Zombie}: exited, awaiting reap (PCB held for exit status); leaks PIDs if unreaped. \textcolor{Bittersweet}{Orphan}: parent died first; reparented to \texttt{init} (PID 1).

\subsection{Context switch}
Save running process's CPU registers + PC + MMU state into its PCB; restore next process's. Triggered by \textcolor{Bittersweet}{timer interrupt}, blocking I/O, preemption, or voluntary yield. \textcolor{Bittersweet}{Overhead}: register save/restore, TLB flush, cache pollution.

\subsection{Unix process API}
\begin{itemize}
  \item \texttt{pid\_t \textcolor{Bittersweet}{fork}(void)}: clone caller. Returns child PID to parent, $0$ to child, $-1$ on error.
  \item \texttt{int \textcolor{Bittersweet}{exec*}(path, argv[, envp])}: replace image with new program. Same PID; FDs inherited (unless \texttt{FD\_CLOEXEC}); signal handlers reset. Returns only on failure.
  \item \texttt{pid\_t \textcolor{Bittersweet}{wait}(\&status)}: block until \emph{any} child terminates; reap zombie; return its PID.
  \item \texttt{pid\_t \textcolor{Bittersweet}{waitpid}(pid,\&status,opts)}: wait for specific child; \texttt{WNOHANG} for non-blocking.
  \item \texttt{void \textcolor{Bittersweet}{exit}(status)}: terminate; flush buffers; signal parent. Lower 8 bits of \texttt{status} returned.
  \item \texttt{pid\_t \textcolor{Bittersweet}{getpid}()/\textcolor{Bittersweet}{getppid}()}: own / parent PID.
\end{itemize}

\paragraph{Status macros} (interpret \texttt{status} from \texttt{wait}):
\begin{itemize}
  \item \texttt{\textcolor{Bittersweet}{WIFEXITED}(s)}: true if child exited normally.
  \item \texttt{\textcolor{Bittersweet}{WEXITSTATUS}(s)}: exit code (use only if \texttt{WIFEXITED}).
  \item \texttt{\textcolor{Bittersweet}{WIFSIGNALED}(s)}: terminated by signal.
  \item \texttt{\textcolor{Bittersweet}{WTERMSIG}(s)}: terminating signal number.
\end{itemize}

\paragraph{Copy-on-write (COW) and FD inheritance}
\textcolor{ForestGreen}{Copy-on-write}: \texttt{fork} initially shares parent's pages read-only; on first write, kernel page-faults and copies. Avoids full memory copy.
\textcolor{ForestGreen}{FD inheritance}: child inherits parent's open FDs, sharing the same file-offset entry; both see seeks and writes.

\paragraph{fork() process tree}
$n$ \texttt{fork()}s $\Rightarrow 2^n$ procs; \texttt{printf} runs $2^k$ times after $k$ forks. \texttt{fork() \&\& fork()} $\Rightarrow$ 3 (child short-circuits on 0).

\paragraph{Canonical fork-exec-wait}
\begin{verbatim}
pid_t pid = fork();
if (pid == -1)     { /* error */ }
else if (pid == 0) {        /* child */
    execv("/bin/prog", argv);
    _exit(1);  /* exec failed */
} else {                   /* parent */
    int s; waitpid(pid, &s, 0);
    if (WIFEXITED(s))
        /* WEXITSTATUS(s) */;
}
\end{verbatim}

\paragraph{exec* family}
\begin{itemize}
  \item \texttt{l} = list args (variadic, NULL-terminated); \texttt{v} = vector \texttt{argv[]}.
  \item \texttt{p} = search \texttt{PATH} on bare names; \texttt{e} = pass custom \texttt{envp[]}.
  \item Combinations: \texttt{execl}, \texttt{execv}, \texttt{execlp}, \texttt{execvp}, \texttt{execle}, \texttt{execve}. Only \texttt{execve} is the actual syscall; others are libc wrappers.
\end{itemize}
